%======================================================================
% Section 4: Classical and Generalized Thermoelastic Theories
%======================================================================
\section{Classical and Generalized Thermoelastic Theories}
\label{sec:thermoelastic}

This section presents the historical development, governing equations, wave propagation characteristics, and dissipation structure of classical and generalized thermoelastic theories. Each theory is introduced with its physical motivation, mathematical formulation, and domain of applicability. Representative equations are selected to illustrate each theory family; derivations from first principles are given where they illuminate the theory's structure. Citations to primary sources are provided throughout to distinguish established results from open questions.

\subsection{Historical evolution}
\label{sec:4.1}

The development of generalized thermoelasticity spans nearly two centuries, from Fourier's heat equation to modern memory-dependent and nonlocal formulations. Table \ref{tab:T1-historical} summarizes the principal milestones.

Historical context only: the review is SECONDARY-SOURCE, not equation-level evidence (BACKGROUND-ONLY) \cite{Chandrasekharaiah1998}.

\begin{table}[htbp]
\centering
\small
\caption{Selected publication milestones and broad modelling mechanisms. Dates and bibliographic matches are inherited metadata; physical and equation-level claims remain subject to primary-source review.}
\label{tab:T1-historical}
\begin{tabularx}{\textwidth}{@{}p{3.0cm}c p{2.4cm}X X@{}}
\toprule
\textbf{Model/family} & \textbf{Year} & \textbf{Source lead} & \textbf{Broad modelling feature} & \textbf{v1.1 evidence status} \\
\midrule
Classical coupled thermoelasticity & 1956 & \cite{Biot1956} & Coupled mechanical and thermal fields with Fourier conduction & Metadata inherited; equations not re-audited here \\
Lord--Shulman & 1967 & \cite{Lord1967} & Heat-flux relaxation model & Metadata inherited; constitutive details not full-text-checked \\
Green--Lindsay & 1972 & \cite{Green1972} & Single temperature; rate dependence with two relaxation times & Description follows the corrected manuscript; primary-source check open \\
Green--Naghdi II/III & 1991/1993 & \cite{Green1991,Green1993} & Thermal-displacement formulations with distinct type labels & Type-specific claims not independently checked \\
Dual/triple phase lag & 1995/2007 & \cite{Tzou1995,RoyChoudhuri2007} & DPL and Roy Choudhuri first-order TPL forms & Exact expansions and admissibility not checked \\
Eringen-type nonlocality & 1972 onward & \cite{Eringen1972,Eringen1974,Eringen1983} & Spatially nonlocal constitutive response & Kernel and boundary formulation require source review \\
Memory-dependent operator & 2011 onward & \cite{Wang2011} & History-dependent heat-transport formulation & Operator definition and limiting behavior not checked \\
Moore--Gibson--Thompson & 1960/1972; thermoelasticity 2019 & \cite{MooreGibson1960,Quintanilla2019} & Relaxation-origin equation and later thermoelastic formulation; Thompson (1972) attribution not source-verified here & Lineage metadata checked in part; Thompson reference remains UNVERIFIED \\
\bottomrule
\multicolumn{5}{@{}p{\textwidth}@{}}{\footnotesize This table is a navigation aid, not a validated chronology, taxonomy, or comparison of propagation speeds and dissipation.}\\
\end{tabularx}
\end{table}


For the historical overview, see Figure~\ref{fig:theory_evolution} in Section 1.

A recurring motivation for some generalized heat-transport models is that Fourier conduction is parabolic and has formally infinite propagation speed. Generalized constitutive choices explore alternatives for selected regimes, but finite-speed behavior and thermodynamic consistency are model-specific properties that require derivation; neither follows from the label “generalized.”

Joseph and Preziosi's review (1989) and 1990 addendum summarize the second-sound evidence and the objections raised in the intervening discussion \cite{Joseph1989,Joseph1990}. This debate helped motivate later finite-speed heat-transport models; it does not establish a universal need for them.

\subsection{Classical coupled thermoelasticity (Fourier heat conduction)}
\label{sec:4.2}

The classical coupled thermoelasticity theory, developed by Biot (1956) \cite{Biot1956}, combines the equations of linear elasticity with Fourier's law of heat conduction:

\begin{equation}
\rho\ddot{u}_i = (\lambda + \mu)u_{j,ji} + \mu u_{i,jj} - \beta T_{,i},
\label{eq:4.1}
\end{equation}

\begin{equation}
\rho c_v \dot{T} = k\nabla^2 T - \beta T_0 \dot{u}_{i,i},
\label{eq:4.2}
\end{equation}

where $\rho$ is density, $\lambda$ and $\mu$ are Lam\'e constants, $\beta$ is the thermal stress coefficient, $c_v$ is specific heat at constant strain, $k$ is thermal conductivity, and $T_0$ is the reference temperature.

\subsubsection{Wave speeds}

The classical theory predicts two elastic wave types:

\begin{equation}
c_p = \sqrt{\frac{\lambda + 2\mu}{\rho}},
\qquad
c_s = \sqrt{\frac{\mu}{\rho}}.
\label{eq:4.3}
\end{equation}

The thermal field satisfies a parabolic diffusion equation, predicting infinite thermal propagation speed --- the classical paradox that motivated all subsequent generalized theories.

\subsubsection{Physical interpretation}

Equations \eqref{eq:4.1}--\eqref{eq:4.2} are two-way coupled: volumetric strain rate sources temperature, while the temperature gradient contributes to mechanical stress. In the fully dynamic problem neither field can in general be solved independently; sequential treatments are approximations requiring explicit scale assumptions.

\subsection{Lord--Shulman (LS) theory (1967)}
\label{sec:4.3}

Lord and Shulman (1967) \cite{Lord1967} introduced a single thermal relaxation time $\tau_0$ into Fourier's law:

\begin{equation}
q_i + \tau_0 \dot{q}_i = -k T_{,i},
\label{eq:4.4}
\end{equation}

where $\tau_0 > 0$ is the relaxation time and $q_i$ is the heat flux vector.

\subsubsection{Modified heat equation}

Combining \eqref{eq:4.4} with the energy balance yields the LS heat equation:

\begin{equation}
\rho c_v \left(\dot{T} + \tau_0 \ddot{T}\right)
= k\nabla^2 T
- \beta T_0 (\dot{u}_{i,i} + \tau_0 \ddot{u}_{i,i}).
\label{eq:4.5}
\end{equation}

\subsubsection{Finite thermal wave speed}

The thermal wave speed in LS theory is:

\begin{equation}
c_T = \sqrt{\frac{k}{\rho c_v \tau_0}}.
\label{eq:4.6}
\end{equation}

For the displayed LS/Cattaneo law, the formal thermal characteristic speed is finite when $\tau_0>0$. The limiting substitution $\tau_0\to0$ recovers Fourier conduction at the level of the constitutive equation; convergence of solutions is a separate question. These statements are conditional on the displayed model and have not been independently checked against the Lord--Shulman primary text in this revision.

\subsubsection{Well-posedness}

For the displayed Cattaneo-type heat-flux law, $\tau_0>0$ gives a finite characteristic speed for the thermal branch. This parameter condition alone is not a proof that the full coupled initial-boundary-value problem is well-posed; the latter depends on coefficients, domain, and boundary/initial data. A separate fractional Lord--Shulman formulation couples a Klein--Gordon nonlocal operator to the heat law; its printed $\alpha\to0$ branch is not used here as a limiting identity (CONDITIONAL) \cite{Das2026}.

\subsection{Green--Lindsay (GL) theory (1972)}
\label{sec:4.4}

Green and Lindsay (1972) introduced a temperature-rate-dependent thermoelastic theory with two relaxation times and a single temperature field. It is distinct from two-temperature thermoelasticity. The original formulation retains Fourier heat flux, $q_i=-kT_{,i}$, with no flux-relaxation time. A separate rate-dependent two-temperature relation uses its own assumption $t_1\geq t_2>0$ and likewise carries no flux-relaxation term; it is not part of the original single-temperature GL law (CONDITIONAL) \cite{Srivastava2026}.

\subsubsection{Governing equations}

\begin{equation}
\rho c_E (\dot{T} + \tau_0\ddot{T}) = k\nabla^2 T - \beta T_0\dot{u}_{i,i},
\label{eq:4.9}
\end{equation}

The original Green--Lindsay 1972 energy equation has no $\tau_0\ddot{u}_{i,i}$ term. The LS equation does carry $\tau_0$ on the strain-rate coupling; applying the same operator to that term while retaining the GL stress law produces an LS/GL hybrid, not original GL. Later modified GL models may use such couplings; Yu, Xue, and Tian (2018) explicitly present a modified variant \cite{Yu2018ModifiedGL}. The control-suggested Bera et al. item remains unidentified. The cited ``Singh (2020), Symmetry 12:488'' identifier resolves to a Saeed--Abbas--Marin paper, so that author attribution and its relevance to this variant are not asserted here pending the Crossref audit. The condition $\tau_1\geq\tau_0>0$ is primary-source-checked from Green and Lindsay's entropy-production restriction \cite{Green1972}.

\begin{equation}
\sigma_{ij} = \lambda\varepsilon_{kk}\delta_{ij} + 2\mu\varepsilon_{ij} - \beta(T + \tau_1\dot{T})\delta_{ij}.
\label{eq:4.10}
\end{equation}

\subsubsection{Thermal-rate coupling and admissibility}

The second relaxation time enters the stress through the temperature rate; it does not introduce a second temperature field. Green--Lindsay theory is distinct from Chen--Gurtin two-temperature thermoelasticity. The displayed ordering applies to the original GL convention; it is not a claim that every later GL-labelled variant has the same admissible range.

\subsection{Green--Naghdi type II theory (1993)}
\label{sec:4.5}

The type-II system treated here is Green and Naghdi's 1993 thermoelasticity without energy dissipation. Let $\alpha_{\mathrm{GN}}$ be the thermal displacement, with $\dot{\alpha}_{\mathrm{GN}}=T$; it is distinct from the fractional order $\alpha$.

\subsubsection{Thermal displacement gradient formulation}

For the source-free form, time integration with zero integration constant gives the source-checked relation \cite{Green1993}:

\begin{equation}
\kappa^*\nabla^2\alpha_{\mathrm{GN}} = \rho c_e \ddot{\alpha}_{\mathrm{GN}} + \beta T_0 \dot{u}_{i,i}.
\label{eq:4.11}
\end{equation}

The strain-rate coupling carries one fewer time derivative than the inherited form. Here $\kappa^*$ is the thermal-displacement coefficient; below, $\kappa$ denotes the coefficient on the temperature-rate term in GN-III. In the displayed isotropic convention, the candidate entropy flux uses the same scalar coefficient:

\begin{equation}
p_i = \kappa^*\alpha_{\mathrm{GN},i},
\label{eq:4.12}
\end{equation}

where $p_i$ is the entropy-flux vector. No anisotropic tensor generalization is asserted here.

\subsubsection{Energy conservation}

The 1993 type-II paper states that the displayed system permits harmonic-wave propagation without damping. This source-checked property is limited to that system and its stated coefficients; the schematic balance below is not a derived GN-II energy identity.

\begin{equation}
\frac{\partial\mathcal{E}}{\partial t} + \partial_i J_i = \mathcal{P}.
\label{eq:4.13}
\end{equation}

Wave attenuation and energy conservation for type II must be obtained from the specific model equations and boundary-value problem; no application suitability claim is made here. A separate rheological network derivation relates CV/GN/MCV/MGN flux laws and includes second-order coefficients proportional to $\tau^2$; those coefficients do not establish energy ratios or admissibility for the displayed system (CONDITIONAL) \cite{Huang2025}.

\subsection{Green--Naghdi type III theory (1992)}
\label{sec:4.6}

Green and Naghdi's 1992 type-III system contains both the temperature-rate and thermal-displacement contributions. The displayed source-free isotropic equation uses $\kappa$ for the rate term and $\kappa^*$ for the non-dissipative term, consistently with Eq.~\eqref{eq:4.11} \cite{GreenGN1992}.

\subsubsection{Governing equations}

\begin{equation}
\rho c_e\ddot{T} + \beta T_0\ddot{u}_{i,i} = \kappa\nabla^2\dot{T} + \kappa^*\nabla^2 T.
\label{eq:4.14}
\end{equation}

A distinct nonlocal GN-III formulation prints a flux law containing both the temperature-gradient and thermal-displacement-gradient contributions; it is not substituted into the local source-free form shown here (CONDITIONAL) \cite{Biswas2020a}. This source-free form is conditional on the displayed linear isotropic system; the source term and material normalization must be restored for other conventions. The transcribed restrictions include $\kappa\geq0$, $c\geq0$, and $\kappa^{**}=0$; no additional sign restriction on $\kappa^*$ is inferred here.

Any proposed type-III-to-type-II or classical limit must be verified by substitution into the exact system and by a convergence analysis. No general limit theorem is asserted here.

\subsection{Dual-phase-lag (DPL) theory (Tzou, 1995)}
\label{sec:4.7}

Tzou (1995) introduced two phase lags into Fourier's law:

\begin{equation}
q_i(t + \tau_q) = -k T_{,i}(t + \tau_T),
\label{eq:4.15}
\end{equation}

where $\tau_q$ is the heat flux phase lag and $\tau_T$ is the temperature gradient phase lag.

\subsubsection{Taylor expansion}

First-order Taylor expansion about $t$:

\begin{equation}
q_i + \tau_q\dot{q}_i = -k(T_{,i} + \tau_T\dot{T}_{,i}).
\label{eq:4.16}
\end{equation}

Some later DPL variants retain a $\tau_q^2/2$ term in a second-order flux expansion (Yang et al.; Gupta et al.); these are distinct model-specific truncations, not part of Eq.~\eqref{eq:4.16}. No universal lag-order admissibility follows (CONDITIONAL) \cite{Yang2021,Gupta2023}. Formal reductions can sometimes be obtained by setting one or both lag parameters to zero in a specified DPL equation, but the resulting equations and solution limits must be checked for that model. The ordering $\tau_q\geq\tau_T$ is not asserted here as a universal thermodynamic-admissibility condition; a model-specific free-energy/entropy derivation is required.

\paragraph{Guyer--Krumhansl (GK) phonon transport.}

Guyer and Krumhansl's 1966 kinetic treatment derives the nonlocal phonon-flux structure from the phonon Boltzmann equation, not from a ready-made continuum thermoelastic law. In the rapid-normal-scattering regime ($\Omega\tau_N\ll1$), its characteristic length is $\lambda^2=c^2\tau_N\tau_z/5$, where $\tau_z$ is the Ziman-averaged resistive-process relaxation time \cite{GuyerKrumhansl1966}. As a SECONDARY continuum statement, Tzou (1995, Eq.~15) prints $\dot{\bm q}+(c^2C_p/3)\nabla T+\bm q/\tau_R=(\tau_Nc^2/5)[\nabla^2\bm q+2\nabla(\nabla\cdot\bm q)]$; this particular reduction is not attributed to the 1966 paper. Tzou's Eq.~17 maps this phonon case to DPL parameters $\alpha=\tau_Rc^2/3$, $\tau_T=9\tau_N/5$, and $\tau_q=\tau_R$ \cite{Tzou1995}. These continuum and mapping statements are conditional on Tzou's stated phonon approximation; they are not universal DPL identities.

\subsection{Triple-phase-lag (TPL) theory (Roy Choudhuri, 2007)}
\label{sec:4.8}

Roy Choudhuri (2007) \cite{RoyChoudhuri2007} proposed the three-lag TPL law:

\begin{equation}
\begin{aligned}
q_i(t+\tau_q) &= -kT_{,i}(t+\tau_T)\\
&\quad - k^*v_{,i}(t+\tau_v), \qquad \dot v=T.
\end{aligned}
\label{eq:4.17}
\end{equation}

The cited source gives a first-order expansion in the three lags with $0\leq\tau_v<\tau_T<\tau_q$. We use no $\tau_\nu/2$ term. Second-order DPL variants retain a $\tau_q^2/2$ term, but they are distinct from this first-order Roy Choudhuri law and establish no shared admissibility condition (CONDITIONAL) \cite{Yang2021,Gupta2023}. Other published variants include second-order DPL/TPL truncations with a free coefficient or tensor conductivity \cite{Kumar2023,Singh2019,Biswas2019}; a higher-order two-temperature TPL model uses a temperature lag rather than $\tau_v$ \cite{Abouelregal2020}. Further formulations combine memory-kernel bounds \cite{Singh2020a} or nonlocal operators on both stress and inertia \cite{Bayones2021}; related DPL/TPL model studies are also reported by Tibullo et al. \cite{Tibullo2017}. These are distinct constitutive variants, not one universal TPL law; no common stability or admissibility condition is inferred (CONDITIONAL).

\subsection{Moore--Gibson--Thompson (MGT) theory}
\label{sec:4.9}

Moore and Gibson (1960) derived a third-order acoustic wave equation for gas dynamics with relaxation; it is not a heat-conduction law for thermoelastic solids. Later Moore--Gibson--Thompson thermoelastic models are analogues that borrow this higher-order structure \cite{MooreGibson1960,Quintanilla2019}. A candidate constitutive system used in some later thermal formulations is

\begin{equation}
\tau\dot{q}_i+q_i=-kT_{,i}-k^*\vartheta_{m,i},\qquad \dot{\vartheta}_m=T.
\label{eq:4.18}
\end{equation}

Here $\vartheta_m$ is a thermal-displacement variable and $k^*$ must have units consistent with the selected model. This constitutive variant is discussed below with its source-specific assumptions and stability scope.

% R3D-BEGIN:MGT
\paragraph{Thermoelastic MGT evidence.}

The displayed Eq.~\eqref{eq:4.18} is the thermoelastic MGT variant with relaxation in the heat-flux law and no relaxation term in the stress. Quintanilla's Eqs.~(2.4), (2.5), and (3.5) give the source-specific coercivity condition $K_{ij}=k_{ij}-\tau k^*_{ij}$ positive definite, with the scalar condition $k-\tau k^*>0$; in one dimension, $k>k^*\tau$ yields exponential stability, while $k<k^*\tau$ gives instability (PRIMARY-SOURCE-CHECKED) \cite{Quintanilla2019}. The source uses $\alpha_{,j}$ in (2.1) but $\theta_{,j}$ in (2.5) and (3.5), without reconciling that notation; the conclusion is kept model-specific.

Bazarra et al. document a distinct radial thermoelastic MGT variant in which $\tau$ also appears in the stress coupling, $-\beta(\theta_{,i}+\tau\dot{\theta}_{,i})$; their radial decay result assumes $\bar k=k-\tau k^*>0$ (PRIMARY-SOURCE-CHECKED) \cite{Bazarra2021MGTradial}. Thus the relaxation placement in the stress law is variant-specific.

\paragraph{Acoustic MGT evidence.} Acoustic MGT-with-memory PDE decay uses its own acoustic parameter set and is separate from thermoelastic admissibility. The gas-dynamics parameter $\gamma_{\rm ac}=b-\tau c^2$ is acoustic only; it is not the thermoelastic condition $k-\tau k^*>0$.

For the separate abstract acoustic MGT equation with memory, Lasiecka and Wang report exponential decay in the subcritical regime $\alpha\beta>\gamma$ (PRIMARY-SOURCE-CHECKED, ABSTRACT-LEVEL) \cite{LasieckaWang2015MGT}.

Dell'Oro, Lasiecka, and Pata state that at $\alpha\beta=\gamma$ exponential stability holds if and only if $A$ is bounded; otherwise the energy tends to zero without an exponential rate (PRIMARY-SOURCE-CHECKED, ABSTRACT-LEVEL) \cite{DellOroLasieckaPata2016MGT}.

% R3F-BEGIN:ACOUSTIC-MGT
Marchand, McDevitt, and Triggiani analyze an acoustic Jordan--MGT equation and report an exponential-decay result for that acoustic system; this source is not thermoelastic MGT evidence (PRIMARY-SOURCE-CHECKED; ACOUSTIC) \cite{MarchandMcDevittTriggiani2012}.
% R3F-END:ACOUSTIC-MGT

\paragraph{MGT evidence scope.} Limited primary sources are available for the s.4.9--s.4.10 MGT subtopics; these model-specific results do not constitute a family-wide classification.
% R3D-END:MGT

\subsection{Memory-dependent derivative (MDD) models (\texorpdfstring{Wang \& Li}{Wang and Li}, 2011)}
\label{sec:4.10}

MDD models replace the classical time derivative with an integral-form memory operator:

\begin{equation}
D_{\omega_M} f(t) = \frac{1}{\omega_M}\int_{t-\omega_M}^{t} \kappa(t-\xi) f'(\xi)\,d\xi,
\label{eq:4.19}
\end{equation}

where $\omega_M>0$ is the finite horizon in this displayed candidate operator; the subscript distinguishes it from angular frequency elsewhere in the review. Its kernel bounds, normalization, regularity, and any admissible parameter set require verification for the exact formulation. A fading-memory condition on an unbounded lag variable cannot be inferred from this finite-horizon formula. A limit as $\omega_M\to0$ may be considered only after kernel normalization and regularity assumptions are specified; no classical-derivative limit or material calibration is asserted here. Li et al. use a $1/\omega$ normalization and report a degenerate-limit equivalence between their memory-dependent and fractional models; that result is formulation-specific and does not establish a limit for the candidate operator displayed here (CONDITIONAL) \cite{Li2019}.

\subsection{Fractional thermoelasticity}
\label{sec:4.11}

Fractional thermoelasticity generalizes the heat equation using fractional derivatives:

\begin{equation}
{}^C D_t^\alpha T = a_\alpha\nabla^2 T, \qquad 0 < \alpha \leq 2,
\label{eq:4.20}
\end{equation}

where ${}^C D_t^\alpha$ is the Caputo fractional derivative of order $\alpha$ and $a_\alpha$ has dimensions $\mathrm{m^2\,s^{-\alpha}}$ in this scalar equation. Thus $a_1$ is a thermal diffusivity, while $\sqrt{a_2}$ has dimensions of speed in the special case $\alpha=2$. For $1<\alpha\leq2$ this is a fractional diffusion--wave equation; that terminology alone does not establish a thermodynamically admissible, causal, or well-posed coupled thermoelastic model. Initial data and a constitutive embedding must be specified for a physical application.

\paragraph{R3c fractional-source scope.} Primary-source coverage is limited: the sole s.4.11 plan row is metadata-only and not technically eligible. GAP-12 and GAP-13 remain open; no model-specific admissibility claim is inferred from that record.

\subsection{Comparative summary}
\label{sec:4.12}

Table \ref{tab:T2-generalized} summarizes selected constitutive distinctions and the evidence limits for the model classes discussed here; it does not assert family-wide wave-speed, dissipation, or validity results.

\begin{table}[htbp]
\centering
\small
\caption{Selected generalized thermoelastic model classes. The table records broad constitutive distinctions; propagation, admissibility, and energy-flux claims require model-specific primary-source checks.}
\label{tab:T2-generalized}
\begin{tabularx}{\textwidth}{@{}p{2.4cm}XXX@{}}
\toprule
\textbf{Model class} & \textbf{Broad constitutive feature} & \textbf{Status in this revision} & \textbf{Source lead} \\
\midrule
Classical coupled thermoelasticity & Fourier heat flux coupled to mechanical balance & Classical local flux derived; equation set not wholly re-audited & \cite{Biot1956} \\
Lord--Shulman & Heat-flux relaxation & Standard form shown; primary-source equation check pending & \cite{Lord1967} \\
Green--Lindsay & Temperature-rate-dependent, single-temperature formulation with two relaxation parameters & Corrected description; full equations and admissibility unverified & \cite{Green1972} \\
Green--Naghdi type II & Thermal displacement; nondissipative system as stated in the 1993 model & Displayed source-free Eq. (4.11); energy identity remains open & \cite{Green1993} \\
Green--Naghdi type III & Coupled temperature-rate and thermal-displacement terms & Displayed Eq. (4.14); dissipation depends on coefficients and the selected system & \cite{GreenGN1992} \\
Dual/triple phase lag & Lagged constitutive formulations & Exact expansions and parameter restrictions unverified & \cite{Tzou1995,RoyChoudhuri2007} \\
Moore--Gibson--Thompson-type models & Later thermal analogues of a third-order relaxing-gas acoustic equation; not Moore--Gibson 1960 heat conduction & Variant-specific equations and stability open & \cite{MooreGibson1960,Quintanilla2019} \\
Memory-dependent & Kernel-based/history-dependent heat transport & Operator and limiting behavior not full-text-checked & \cite{Wang2011} \\
Fractional scalar diffusion--wave example & Caputo time derivative in the displayed scalar equation; not by itself a coupled thermoelastic model & Order-dependent coefficient units stated; embedding, initial data, and admissibility open & \cite{Povstenko2004,Povstenko2009} \\
\bottomrule
\multicolumn{4}{@{}p{\textwidth}@{}}{\footnotesize No universal finite-speed, dissipation, stability, or well-posedness result is inferred from a family label.}\\
\multicolumn{4}{@{}p{\textwidth}@{}}{\footnotesize Generalized energy-flux identities remain OPEN unless individually derived from the selected model.}\\
\end{tabularx}
\end{table}


Published examples include a hyperbolic two-temperature closure \cite{Prasad2020}, fractional-order heat conduction with two temperatures \cite{Ezzat2011}, semiconductor applications of a hyperbolic two-temperature model \cite{Alzahrani2020}, and a fractional two-temperature hybrid with relaxation times \cite{Bassiouny2020}. These examples illustrate distinct constructions only; they do not establish shared admissibility or stability results (CONDITIONAL).

\begin{figure}[htbp]
\centering
\includegraphics[width=0.95\textwidth]{fig07-theory-map.pdf}
\caption{Schematic map of selected thermoelastic model families. Parameter limits and links are model-specific; no universal reduction or convergence result is asserted. Green--Lindsay is represented as a single-temperature, temperature-rate-dependent theory, distinct from two-temperature models.}
\label{fig:theory_map}
\end{figure}

\subsection{State variables and constitutive distinctions}
\label{sec:4.13}

The model classes surveyed in this section differ in their chosen state variables and constitutive operators. Relaxation models may promote heat flux to an evolving variable; two-temperature formulations distinguish temperature-like fields; thermal-displacement formulations introduce a variable whose time derivative is related to temperature; phase-lag models shift the arguments of constitutive quantities; hereditary and fractional models use history operators. These descriptions are schematic. Their thermodynamic interpretation, initial data, modal consequences, and limiting behavior must be derived from each model's exact equations. In particular, Green--Lindsay is not a two-temperature construction.

\subsection{Thermodynamic admissibility as a selection principle}
\label{sec:4.14}

Thermodynamic admissibility must be established separately for each constitutive family by specifying its free energy, entropy, state variables, evolution laws, and admissible processes. The inherited universal dissipation claims and parameter exclusions have not been re-derived here. No parameter region is excluded, and no Green--Lindsay admissibility result is asserted in this revision without a primary-source derivation.

Thermodynamic restrictions must be derived from the free energy, entropy, heat-flux law, and internal variables of each chosen model. No single dissipation expression or parameter-ordering condition applies to all relaxation, two-temperature, phase-lag, memory-dependent, or fractional formulations. The family-wide assertions in the inherited text are not supported by derivations in this revision and are withdrawn pending model-specific primary-source and thermodynamic checks.

Parameter restrictions are model-specific and must be included when a model is used. The limits and admissible ranges summarized in inherited tables are not treated as verified unless a derivation or full-text source check is identified.

\subsection{Modal structure is model-specific}
\label{sec:4.15}

The modal structure must be obtained by substituting the chosen constitutive equations into the field equations and solving the resulting characteristic problem. A model label alone does not determine the number of propagating, evanescent, or attenuating roots, nor their physical admissibility. The inherited wave-speed, attenuation, and mode-count claims for the relaxation, two-temperature, thermal-displacement, phase-lag, memory-dependent, and fractional families have not been checked in this revision. Root classification and interface matching are therefore treated as problem-specific tasks in Sections~\ref{sec:methods} and~\ref{sec:interfaces}; no universal count or crossover frequency is asserted here.

\subsection{Notation summary for Section 4}

Table \ref{tab:notation-sec4} summarizes the notation introduced in this section.

\begin{table}[htbp]
\centering
\small
\caption{Notation introduced in Section 4.}
\label{tab:notation-sec4}
\begin{tabularx}{\textwidth}{@{}p{1.8cm}X p{2.2cm}p{2.1cm}@{}}
\toprule
\textbf{Symbol} & \textbf{Meaning} & \textbf{Units} & \textbf{Equation} \\
\midrule
$\tau_0$ & LS relaxation time & s & \eqref{eq:4.4} \\
$\tau_q, \tau_T$ & DPL lag parameters in the displayed inherited form & s & \eqref{eq:4.15}--\eqref{eq:4.16} \\
$c_T$ & Thermal wave speed & m/s & \eqref{eq:4.6} \\
$\beta$ & Thermal stress coefficient & Pa/K & \eqref{eq:4.1} \\
$c_v$ & Specific heat at constant strain & J/(kg$\cdot$K) & \eqref{eq:4.2} \\
$k$ & Thermal conductivity & W/(m$\cdot$K) & \eqref{eq:4.2} \\
$\alpha_{\mathrm{GN}}$ & Thermal displacement (GN-II; distinct from fractional order $\alpha$) & K$\cdot$s & \eqref{eq:4.11} \\
$\kappa,\kappa^*$ & GN-III rate and thermal-displacement coefficients & model-dependent & \eqref{eq:4.11}, \eqref{eq:4.14} \\
$p_i$ & Entropy flux (GN-II) & W/(m$^2$$\cdot$K) & \eqref{eq:4.12} \\
$D_{\omega_M}$ & MDD operator & 1/s & \eqref{eq:4.19} \\
$\omega_M$ & MDD finite time horizon & s & \eqref{eq:4.19} \\
$\kappa(t-\xi)$ & MDD kernel & dimensionless & \eqref{eq:4.19} \\
${}^C D_t^\alpha$ & Caputo fractional derivative & $t^{-\alpha}$ & \eqref{eq:4.20} \\
$\alpha$ & Fractional derivative order & dimensionless & \eqref{eq:4.20} \\
$a_\alpha$ & Coefficient in displayed scalar fractional equation & $\mathrm{m^2\,s^{-\alpha}}$ & \eqref{eq:4.20} \\
\bottomrule
\end{tabularx}
\end{table}

\subsection{Open questions for Section 4}

\begin{enumerate}[label=\textbf{OQ4.\arabic*}]
\item What is the complete time-domain energy flux for LS theory, including the contribution from the flux-relaxation term? \\ \label{oq:4.1}
\item What heat-flux law and energy balance follow from the original single-temperature Green--Lindsay relaxation-time convention? \\ \label{oq:4.2}
\item What is the relationship between energy flux and entropy flux in GN-II theory? \\ \label{oq:4.3}
\item What is the mixed energy flux for GN-III theory? \\ \label{oq:4.4}
\item How do different MDD kernel choices affect the energy flux and wave propagation? \\ \label{oq:4.5}
\item Can rigorous limit theorems with convergence proofs be established between theory pairs? \\ \label{oq:4.6}
\end{enumerate}

% LocalWords:
